\documentclass[11pt,a4paper]{article}
\usepackage[utf8]{inputenc}
\usepackage[T1]{fontenc}
\usepackage[brazil]{babel}
\usepackage{lmodern}
\usepackage[margin=2.1cm]{geometry}
\usepackage{microtype,booktabs,tabularx,enumitem}
\usepackage[hidelinks]{hyperref}
\setlength{\parindent}{0pt}
\setlength{\parskip}{5pt}
\setlist{nosep,leftmargin=*}
\newcommand{\msg}[1]{\texttt{\detokenize{#1}}}
\title{Chat com ordenação\\\large Relatório -- esboço inicial}
\author{Levi \and Henri \and André \and Gabriel}
\date{29 de outubro de 2026}
\begin{document}
\maketitle
\vspace{-1.4em}

\textbf{Estado.} Este arquivo organiza o relatório pedido no enunciado.
A versão final deverá ter 4 a 6 páginas, com a solução realmente implementada,
as decisões efetivas e os testes executados. A extensão atual é provisória.
Ainda não há aplicação ou resultados de execução.

\section{Objetivo e arquitetura proposta}
Construir um chat cuja mesma sequência de mensagens seja entregue aos clientes
conectados a nós diferentes, com recuperação após a queda de um nó.

\begin{tabularx}{\linewidth}{@{}lX@{}}
\toprule
Parte & Proposta inicial \\
\midrule
Aplicação & Python e FastAPI; interface em HTML, CSS e JavaScript puro. \\
Execução & Três nós equivalentes no Docker Compose, inicialmente em uma máquina. \\
Transporte & WebSocket para clientes; HTTP e JSON entre nós. \\
Descoberta & DNS do Compose e nomes de serviços fornecidos por configuração. \\
Ordenação & Multicast total ordenado por Lamport, com ACK de todos na visão. \\
Coordenação & Heartbeat e Bully para escolher quem reorganiza o grupo. \\
Estado & Histórico e pendências em memória, com cópias em cada nó. \\
\bottomrule
\end{tabularx}

\textbf{Princípios.} Funções pequenas, código explicável pelos quatro integrantes
e dependências justificadas. Começar com uma sala, apelido e texto.
O único extra planejado é um painel com nós, líder, época e relógios.
A demonstração em uma máquina é uma proposta, ainda a alinhar com o grupo.

\section{Hipóteses e objetivos de correção}
O cenário inicial admite uma queda por vez, comunicação funcionando entre os
sobreviventes e nenhuma nova falha durante reorganização ou sincronização.
Os timeouts pressupõem respostas dos nós saudáveis dentro dos limites configurados.

Pretende-se demonstrar: mesma ordem de entrega, ausência de duplicatas no histórico,
preservação das mensagens confirmadas após a falha prevista e sincronização de um
nó que retorna. Uma tentativa sem confirmação pode ter sido entregue.

Partições e falsas suspeitas persistentes ficam fora da garantia inicial.
São necessários pelo menos dois dos três nós para operar. Essa regra não prova
ausência de líderes concorrentes nem permite classificar o sistema como CP.
Reiniciar todos os nós perde o histórico.


\section{Algoritmos e decisões de projeto}
\begin{tabularx}{\linewidth}{@{}lX@{}}
\toprule
Problema & Escolha e justificativa inicial \\
\midrule
Ordem de entrega & Multicast por Lamport, fila ordenada e ACKs de todos na visão.
A ordenação é responsabilidade do protocolo distribuído de mensagens. \\
Detecção de falha & Heartbeat com timeout. O timeout produz suspeita;
a reorganização é uma etapa separada. \\
Coordenação & Bully e época identificam o coordenador da visão.
O líder administra o grupo e não numera as mensagens do chat. \\
Recuperação & Pausa, reconciliação de estado e instalação confirmada da visão.
A segurança dessa adaptação didática ainda precisa ser examinada. \\
Reentrada & Snapshot antes de admitir o nó e permitir sua participação na eleição. \\
\bottomrule
\end{tabularx}

O contrato de mensagens e as transições propostas estão em
\msg{docs/protocolo.tex}. O relatório deve explicar escolhas e consequências,
sem reproduzir todos os envelopes do protocolo.

A configuração pequena e o estado em memória reduzem dependências e facilitam
a explicação. Em contrapartida, o histórico depende dos sobreviventes.
A exigência de ACK de todos torna a entrega dependente do membro mais lento;
uma suspeita de falha interrompe o fluxo até a reorganização.

A troca de visão usa pausas locais para coletar histórico e pendências dos
sobreviventes. Cada pausa define uma fronteira de processamento, não um instante
global. O histórico reconciliado deve preservar como prefixo o que cada nó já
entregou. Os detalhes e um exemplo com mensagens em trânsito ficam no protocolo.

\section{Consistência e disponibilidade: análise inicial}
A consistência pretendida é a mesma ordem de mensagens entregues e um histórico
anterior preservado como prefixo após a troca de visão. Isso ainda é um objetivo,
não uma propriedade demonstrada da implementação.

Durante eleição, instalação de visão e sincronização, a proposta pausa a aplicação.
Com um único sobrevivente, novas mensagens ficam indisponíveis.
A opção prioriza preservar o estado conhecido dentro do cenário de falha suportado.

A regra de dois participantes não prova exclusividade de liderança ou consistência
sob partições. A ordenação total proposta também não estabelece, por si só,
linearizabilidade. A análise final de CAP deverá explicitar o modelo de falha,
a semântica de confirmação e o comportamento observado, sem atribuir ao sistema
garantias de partição que não foram implementadas e verificadas.
Como base para essa discussão:
\href{https://groups.csail.mit.edu/tds/papers/Gilbert/Brewer2.pdf}{Gilbert e Lynch,
Perspectives on the CAP Theorem}.

\section{Validação planejada e resultados}
Nenhum cenário abaixo foi executado. Na versão final, registrar condições,
comandos, resultado real e evidência de cada teste selecionado.

\begin{tabularx}{\linewidth}{@{}lX@{}}
\toprule
Cenário planejado & Comportamento esperado, ainda não verificado \\
\midrule
Concorrência & Clientes em nós diferentes recebem a mesma sequência. \\
Retry & Repetir um pacote após resposta perdida não repete seus efeitos. \\
Reconexão & Reenvio com o mesmo UUID não duplica o histórico. \\
Queda do líder & Mensagens confirmadas sobrevivem e os nós restantes retomam. \\
Mensagens em trânsito & A troca de visão preserva o prefixo já entregue. \\
Reentrada & O nó reiniciado sincroniza antes de participar normalmente. \\
Um sobrevivente & O sistema interrompe a aceitação de novas mensagens. \\
\bottomrule
\end{tabularx}

\textbf{Resultados atuais.} Implementação e validação em execução pendentes.
Compilar os documentos não valida os algoritmos ou a recuperação.

\textbf{Revisão do esboço.} Foi examinado um exemplo com queda do nó 3:
o nó 1 entregou uma mensagem que ainda era pendência do nó 2, e uma segunda
mensagem estava em trânsito para o nó 2. Um modelo de dados em memória conferiu
a união dos relatos, a preservação de prefixos, os sufixos a notificar e o
bloqueio diante de um prefixo incompatível. Não houve execução dos nós, da rede
ou do protocolo; essa conferência não demonstra sua correção geral.

\section{Limitações e pontos a consolidar}
Ainda é necessário revisar o bootstrap, os formatos exatos, a confirmação ao
cliente, a deduplicação entre nós, a fronteira de congelamento e a reconciliação.
Partições, falsas suspeitas persistentes e novas falhas durante reorganização
ficam fora da garantia inicial. As limitações efetivas devem ser reavaliadas
conforme a implementação e os experimentos.

\section{Uso de IA}
Até esta versão, ferramentas de IA auxiliaram o planejamento, a revisão das
hipóteses do protocolo e a elaboração dos esboços de documentação.
O plano foi discutido com Codex e revisado com apoio de Claude.
Não há implementação de aplicação a atribuir nesta etapa.
Este registro deve ser atualizado com os usos posteriores e a revisão humana
realizada pelo grupo. Os integrantes precisam compreender e explicar a solução.

\section*{Referências iniciais}
Enunciado do professor: Trabalho Final -- Sistemas Distribuídos 2026,
fornecido durante o planejamento.
\href{https://lamport.azurewebsites.net/pubs/time-clocks.pdf}{Lamport (1978)};
\href{https://homepage.cs.uiowa.edu/~ghosh/Bully.pdf}{Garcia-Molina (1982)};
\href{https://www.cs.cornell.edu/home/rvr/papers/podc/node7.html}{Van Renesse et al.,
comunicação e mudança de visão}.
\end{document}
